% CANDIDATE model note. Not a theorem. Does not edit main.tex.
% XeLaTeX. Sources cited below were opened on 2026-10-04.
\documentclass[11pt]{article}
\usepackage{fontspec}
\usepackage{amsmath,amssymb}
\usepackage[margin=1in]{geometry}
\usepackage{booktabs}
\usepackage{hyperref}
\setmainfont{Latin Modern Roman}
\hypersetup{colorlinks=true,urlcolor=blue,linkcolor=blue}

\title{CANDIDATE: a proof-space limit model\\
\large PROOFSPACE1, not a theorem}
\author{Grok Bot, identity \texttt{j15marti}\\
j15marti/pnp, branch \texttt{main}}
\date{4 October 2026, about 11:49 AM PT (idea); note written the same day}

\begin{document}
\maketitle

P versus NP remains open.

This note defines one \textbf{CANDIDATE} object: a resource space, and the boundary where two of its regimes meet. It is not a theorem, not a separation, and not a claim that any limit of these spaces solves P versus NP. The limit idea does not by itself produce a separation. No paper, year, theorem, or number below is used unless the source was opened in this pass. Items marked ``not verified this pass'' are omitted from the bibliographic claims.

\section*{Standing project facts, not revived}
These are recorded so the model does not walk backward. They are project status, not new computations, and not literature theorems.
\begin{itemize}
\item P versus NP remains open.
\item Two load-bearing arrows remain \textbf{OPEN}: a source-legal characteristic-2 product-span rank, and a compiler from an arbitrary SAT decider into a charged \((T,S,R)\) model.
\item The semantic-cut claim \(\ge 48\) on the negative-star \(D(4,8)\) architecture is \textbf{FALSIFIED/RETIRED}. The exact cut on that architecture is \(8\). The contact count \(48\) stays \textbf{PROVED/EXACT} inside that architecture.
\item \(Q_4/D(4,8)\) is \textbf{FALSIFIED/RETIRED} as a read-2-surviving lower-bound candidate (\(\mathrm{lw}(Q_4)=7\), read-2 size \(\le 64\)).
\item The overnight slice-cover controls \((n,7,7)\) through at least \((311,7,7)\), with \(\tau\)-histogram \(\{1:126\}\), are computational easy-structure only. They are not TASK-001.
\item Contacts are polynomially the same as branching programs. That polynomial agreement does not close the compiler gap.
\end{itemize}

\section*{What was opened}
Opened and read far enough to quote a statement:
\begin{itemize}
\item Albert Atserias and V\'ictor Dalmau, \emph{A Combinatorial Characterization of Resolution Width}, ECCC, title page dated 27 May 2002. PDF:
\url{https://eccc.weizmann.ac.il/report/2002/035/download/}.
\item Eli Ben-Sasson and Avi Wigderson, \emph{Short Proofs Are Narrow---Resolution Made Simple}, \emph{Journal of the ACM} 48(2), March 2001, pages 149--169. PDF:
\url{https://people.inf.ethz.ch/emo/SatSem05/Papers/BensassonWidgerson01.pdf}.
\item Michael Alekhnovich and Alexander A.\ Razborov, \emph{Lower Bounds for Polynomial Calculus: Non-Binomial Case}, manuscript title page dated 13 March 2003. PDF:
\url{http://people.cs.uchicago.edu/~razborov/files/misha.pdf}.
The embedded fonts did not extract as text. The abstract and the first two introduction pages were read from a rendered image of that PDF. Numbered body theorems past those pages were \emph{not} transcribed this pass.
\item Lance Fortnow, \emph{Time-Space Tradeoffs for Satisfiability}. PDF:
\url{https://lance.fortnow.com/papers/files/npvsnl.pdf}.
A separate ECCC download of TR97-041 was saved and did not yield readable text, so no bibliographic detail is taken from that file.
\item Lance Fortnow, Richard Lipton, Dieter van Melkebeek, and Anastasios Viglas, \emph{Time-Space Lower Bounds for Satisfiability}, manuscript dated 7 September 2004. PDF:
\url{https://lance.fortnow.com/papers/files/tst.pdf}.
\item Ryan Williams, \emph{Alternation-Trading Proofs, Linear Programming, and Lower Bounds}, the manuscript whose header identifies the author and the year 2013. PDF:
\url{https://people.csail.mit.edu/rrw/automated-lbs.pdf}.
The exponents \(2\cos(\pi/7)\) and \(n^{1.801}\) are taken from sentences in this opened manuscript. The papers it cites as Williams 2008 and Buss--Williams 2012 were \textbf{not opened this pass}.
\item Paul Beame, Michael Saks, and Jayram S.\ Thathachar, \emph{Time-Space Tradeoffs for Branching Programs}. PDF:
\url{https://homes.cs.washington.edu/~beame/papers/focsbranch.pdf}.
This is the branching-program time/size paper in the Beame--Jayram--Saks line. Wegener and Jukna were not opened.
\end{itemize}

Not opened, so not used: Razborov's 8ECM survey (the URL tried returned 404); any polynomial-calculus \emph{space} paper; Impagliazzo--Pudl\'ak--Sgall beyond mentions inside files that were not themselves the source of a quoted size-degree constant; Wegener; Jukna.

\section*{CANDIDATE, not a theorem}
\textbf{Definition (CANDIDATE).}
A \emph{proof space} is the set of points
\[
\bigl(T,\; W_{\mathrm{BP}},\; S,\; (T_{\mathrm{RAM}},S_{\mathrm{RAM}}),\; d_{\mathrm{PC}},\; w_{\mathrm{Res}},\; \mathrm{CSpace}\bigr)
\]
whose coordinates are, at least, the following.
\begin{itemize}
\item \(T\): time of a correct decider.
\item \(W_{\mathrm{BP}}\): branching-program width in the leveled sense, or, when a source states size and length instead of width, the pair \((\mathrm{length},\mathrm{size})\). In the opened Beame--Saks--Thathachar manuscript, path length is time and the logarithm of the number of nodes is space. This note does not rename that pair as a width theorem.
\item \(S\): mutable work space of a decider.
\item \((T_{\mathrm{RAM}},S_{\mathrm{RAM}})\): simultaneous time and space on a random-access machine, deterministic or nondeterministic as stated by the theorem in force.
\item \(d_{\mathrm{PC}}\): degree of a polynomial-calculus refutation.
\item \(w_{\mathrm{Res}}\): width of a resolution refutation (maximum number of literals in a clause of the proof, as in the opened Ben-Sasson--Wigderson paper).
\item \(\mathrm{CSpace}\): clause space of a self-contained resolution refutation, in the sense of the opened Atserias--Dalmau paper (maximum number of clauses stored in a configuration).
\end{itemize}
A point is \emph{feasible} when some correct decider or some correct proof attains it on the instance in question. A region is \emph{forbidden} when a theorem in an opened source rules that region out for the model that theorem actually talks about. The proposed object is not the whole space. It is the \emph{boundary} at which two regimes meet: clause space against resolution width, resolution width against polynomial-calculus degree, or RAM space against the logarithm of branching-program size. The question the boundary is allowed to ask, and not to answer, is whether the project's \(D(n,m)\) / product-span obstruction sits in a forbidden region, or only in a region that has not been connected to a general machine.

The compiler arrow being OPEN means the second alternative is the present status. Nothing in the limit construction moves a \(D(n,m)\) point into a forbidden region of a general SAT machine.

\section*{What the opened theorems forbid}
\subsection*{Resolution width, clause space, and size}
Atserias--Dalmau, abstract, as read: the minimum space of refuting a 3-CNF is bounded from below by the minimum width of refuting it, minus 3. Section 4, as read: for a \(k\)-CNF, the minimal clause space is bounded from below by the minimal width minus \(k\). Theorem 2 is that comparison. Theorem 1 characterizes resolution width by the Boolean existential pebble game. The additive pebble offset did not survive text extraction cleanly, so it is not restated.

Forbidden, for resolution only: a \(k\)-CNF refutation whose width exceeds its clause space by more than \(k\). Equivalently, the half-space
\[
w_{\mathrm{Res}} > \mathrm{CSpace} + k
\]
contains no correct resolution refutation of that \(k\)-CNF. The boundary \(w_{\mathrm{Res}} = \mathrm{CSpace}+k\) is where clause space meets width up to the shift the theorem allows. Approaching that boundary is a statement about resolution configurations. It is not a statement about arbitrary algorithms.

Ben-Sasson--Wigderson, opened:
\begin{itemize}
\item Theorem 3.3: \(w(\mathcal{F}\vdash 0)\le w(\mathcal{F})+\log S_T(\mathcal{F})\). The proof takes tree-like size \(S_T\le 2^b\) and concludes width at most \(w(\mathcal{F})+b\), so the logarithm is base 2.
\item Corollary 3.4: \(S_T(\mathcal{F})\ge 2^{\,w(\mathcal{F}\vdash 0)-w(\mathcal{F})}\).
\item Theorem 3.5: \(w(\mathcal{F}\vdash 0)\le w(\mathcal{F})+O\bigl(\sqrt{n\ln S(\mathcal{F})}\bigr)\). The proof sets the fat-clause threshold to \(\lceil\sqrt{2n\ln S(\mathcal{F})}\rceil\).
\item Corollary 3.6 is the size consequence of that width bound: general resolution size is exponential in the square of the width gap over \(n\). The displayed formula was extracted imperfectly; the theorem that was read cleanly is Theorem 3.5, and the note relies on that.
\end{itemize}
Forbidden, for resolution only: a refutation that is simultaneously small in the sense of Theorem 3.5 and wider than the right-hand side. The opened Atserias--Dalmau abstract also says that their dense linear order principle has polynomial-size resolution refutations and is a further example that the size-width relationship is tight. So large width does not, by itself, mean large size beyond what Theorem 3.5 already says. Ben-Sasson--Wigderson recall Cook--Reckhow: \(\mathbf{NP}=\mathbf{coNP}\) if and only if some propositional proof system proves every tautology in size polynomial in the tautology. A resolution size lower bound is not that statement.

These theorems do \textbf{not} forbid a polynomial-time algorithm for arbitrary SAT. They do not forbid small proofs in a stronger system than resolution. They do not say that every unsatisfiable formula has large width.

\subsection*{Polynomial calculus degree}
Alekhnovich--Razborov, abstract, as read from the rendered title page: they give a hardness criterion they call immunity, satisfied by a random function, and prove linear lower bounds on the degree of polynomial-calculus refutations for a wide class of tautologies based on immune functions. Applications stated there include \(\mathrm{mod}_p\) Tseitin tautologies in the Boolean case (axioms \(x_i^2=x_i\)), hard for polynomial calculus over fields whose characteristic is different from \(p\); flow tautologies based on majority, hard over any field; and an \(\Omega(n)\) degree lower bound for random \(k\)-CNFs over fields of characteristic 2. The introduction, as read, says binomial methods do not apply once one insists on \(x^2-x=0\) in characteristic 2, and points at random \(k\)-CNF in characteristic 2 as the example.

Forbidden: low polynomial-calculus degree for those formula families, over the fields the abstract names. Not forbidden: polynomial-time SAT; low degree for a formula outside those families; any relation between polynomial-calculus \emph{space} and degree. Polynomial-calculus space was \textbf{not verified this pass}. The coordinate \(d_{\mathrm{PC}}\) therefore has a forbidden region only for the families in that abstract. The coordinate ``polynomial-calculus space'' is named in the candidate only as an axis whose forbidden region is empty in this note.

Width approaching degree is \textbf{not} a theorem of the sources opened here. No simulation between resolution width and polynomial-calculus degree was read. That meeting line is a question inside the candidate, not a forbidden wall.

\subsection*{RAM time and space}
Fortnow, introduction, as read. For every \(\varepsilon>0\),
\begin{align*}
\mathrm{SAT}&\notin \mathbf{NTISP}\bigl[n^{1+o(1)},\, n^{1-\varepsilon}\bigr],\\
\mathrm{SAT}&\notin \mathbf{DTISP}\bigl[n^{1+o(1)},\, n^{1-\varepsilon}\bigr],\\
\mathrm{SAT}&\notin \mathbf{NTIME}\bigl[n^{1+o(1)}\bigr]\cap \mathbf{NL}.
\end{align*}
The same introduction states a more general bound: for \(r(n)\) unbounded with \(r(n)=O(\log n/\log\log n)\) and \(\varepsilon>0\),
\[
\mathrm{SAT}\notin \mathbf{NTIME}\bigl[n^{1+1/r(n)}\bigr]
\cap \mathbf{NTISP}\bigl[n^{o(r(n))},\, n^{1-\varepsilon}\bigr].
\]
It also states, with Buhrman, that SAT has neither log-space uniform \(\mathbf{NC}^1\) circuits of size \(n^{1+o(1)}\) nor log-space uniform branching programs of size \(n^{1+o(1)}\). The model is a Turing machine with random access to the input. The paper says these are time-space tradeoffs, and it discusses why a collapse of a nonconstant level of the polynomial-time hierarchy would be needed before the same diagonalization could separate \(\mathbf{NL}\) from \(\mathbf{NP}\). It does not perform that separation.

Fortnow--Lipton--van Melkebeek--Viglas, Theorem 1.1, as read. Let \(\varphi=(1+\sqrt{5})/2\). For any constant \(c<\varphi\) there is a positive constant \(d\) such that SAT cannot be solved on a deterministic random-access Turing machine in time \(n^c\) and space \(n^d\). As \(c\) approaches 1 from above, \(d\) approaches 1 from below. The introduction says that, as the space bound becomes subpolynomial, the time exponent approaches \(\varphi\approx 1.618\). Theorem 1.2: on conondeterministic random-access Turing machines the same shape holds for every constant \(c<\sqrt{2}\). The introduction also says that, on unrestricted random-access machines, no time lower bound for SAT better than the trivial linear one was known in that manuscript. The bounds apply to nondeterministic linear time and, by the reductions they invoke, to the natural \(\mathbf{NP}\)-complete problems they describe. They are simultaneous time \emph{and} space bounds.

Williams, opened manuscript, as read. The best time lower bound it records for SAT on algorithms using \(n^{o(1)}\) space is
\[
n^{2\cos(\pi/7)-o(1)}\ge n^{1.801},
\]
attributed there to Williams 2008 (that paper was not opened). The same manuscript proves that any SAT algorithm running in time \(t(n)\) and space \(s(n)\) must satisfy \(t(n)\cdot s(n)\ge n^{\varphi-o(1)}\) infinitely often, with \(\varphi\approx 1.618\) the golden ratio, improving a product bound it attributes to Fortnow--Lipton--van Melkebeek--Viglas as \(n^{1.573}\). It also records that Buss and Williams (2012; not opened) showed \(2\cos(\pi/7)\) is optimal for alternation-trading proofs in the \(n^{o(1)}\)-space setting: the method cannot prove an \(n^2\) time lower bound of that form. The manuscript writes \(2\cos(\pi/7)\approx 1.8019\). So \(2\cos(\pi/7)\) is confirmed \emph{as a cited lower bound and as a limit of one proof method}, inside this opened paper. It is not a tight complexity of SAT, and it is not a polynomial-space lower bound.

Forbidden: the simultaneous corners above, on the machine models those theorems name. Not forbidden: SAT in polynomial time and polynomial space; SAT in time \(n^{2}\) and space \(n\); SAT in deterministic time \(n^{c}\) for \(c\ge \varphi\) when the space exponent is allowed to be large enough that the theorem's \(d\) does not apply; any separation of \(\mathbf{P}\) from \(\mathbf{NP}\).

\subsection*{Branching-program length and size}
Beame--Saks--Thathachar, as read. There is an explicit function, the Boolean quadratic form \(BQF_M\) on a Sylvester matrix over \(\mathrm{GF}(3)\), such that every branching program of length at most \(1.0178n\) has size at least \(2^{n/3}\), and therefore exponential size. Theorem 2 states that there is an \(\varepsilon>0\) such that length \((1+\varepsilon)n\) forces exponential size. The explicit check in the same manuscript says that \(\varepsilon\le 0.0178\) gives size at least \(2^{n/3}\). The manuscript also separates syntactic and semantic read-\(k\) branching programs for \(k>1\), by a function with polynomial-size semantic read-twice programs and exponential size for syntactic read-\(k\).

The introduction of that manuscript says that an explicit function in \(\mathbf{P}\) with no linear-length polynomial-size branching program was still open there, because the model is nonuniform and Fortnow's bound does not transfer, and because ruling out length \(1.0178n\) does not rule out every linear length. This note does not upgrade \(1.0178n\) into an \(O(n)\) prohibition.

Forbidden: polynomial size at length \(\le 1.0178n\) for that one function. Not forbidden: short branching programs for SAT; a general width lower bound for arbitrary polynomial-time SAT; the identification of RAM space with the logarithm of branching-program width for a uniform machine. That identification is a proposed boundary of the candidate. It was not verified as a simulation theorem this pass. The only uniform branching-program statement read here is Fortnow's log-space uniform size \(n^{1+o(1)}\) prohibition for SAT.

\section*{Where the project's obstruction sits}
The product-span arrow and the compiler into a charged \((T,S,R)\) model are OPEN. Contacts match branching programs up to a polynomial and still do not compile an arbitrary SAT decider. Therefore a point that lives only in the \(D(n,m)\) contact or product-span geometry is, on this evidence, a point of a restricted architecture. It is not shown to be a forbidden point of resolution, of polynomial calculus, of RAM time-space, or of general branching programs.

The retired claims are not boundaries of this candidate. Semantic cut \(\ge 48\) is false on negative-star \(D(4,8)\) (exact cut \(8\)). Contact count \(48\) remains an exact count inside that architecture, not a general-machine lower bound. \(Q_4\) with \(\mathrm{lw}(Q_4)=7\) and read-2 size at most \(64\) is not a read-2-surviving lower-bound candidate. Slice-cover data through \((311,7,7)\) are easy-structure computations, not a proof-space boundary.

\section*{One finite test, unrun}
No numerical frontier is stated, because none was computed.

\textbf{TEST-PROOFSPACE-D48 (UNRUN).}
Choose one fixed CNF \(F\) already determined by a repository encoding of either the negative-star \(D(4,8)\) decision or one slice-cover instance among those already run through \((311,7,7)\). This note does not choose the clauses and does not run a solver. On that fixed \(F\), compute two numbers and nothing else:
\begin{enumerate}
\item the minimum resolution width of \(F\);
\item the minimum polynomial-calculus degree of \(F\) over a field of characteristic 2.
\end{enumerate}
The test asks only whether those two coordinates meet, differ by a shift comparable to the initial width, or stay apart. If \(F\) is a \(k\)-CNF, a third recorded number is the minimum clause space, to be compared with width minus \(k\) under the Atserias--Dalmau inequality. A further pair, RAM work space of a direct decider against the logarithm of branching-program size for the same predicate, is part of the question and is not a theorem. Status of every number in this test: \textbf{UNRUN}.

\section*{What the limit does not say}
Taking a limit as width approaches degree, as clause space approaches width, or as RAM space approaches the logarithm of branching-program size, produces a boundary in the candidate space. A boundary is not a diagonalization. None of the opened theorems becomes a superpolynomial lower bound for arbitrary polynomial-time SAT by being rewritten as a limit. If a future compiler placed every polynomial-time SAT decider on the feasible side of one of the \emph{verified} forbidden regions above, that would be a new theorem, and it is not this note.

\section*{Barrier audit}
\begin{itemize}
\item \textbf{Relativization.} Not claimed. Fortnow's opened manuscript already records a relativized limit on one proposed route from these time-space ideas to \(\mathbf{NP}\ne\mathbf{NL}\). This candidate does not cross that limit.
\item \textbf{Algebrization.} Not claimed. Not checked against an algebrizing formulation this pass.
\item \textbf{Natural proofs.} Not claimed. The candidate is a map of known resource bounds, not a circuit lower bound.
\end{itemize}

\section*{Status table}
{\footnotesize
\begin{tabular}{@{}p{0.40\textwidth}p{0.22\textwidth}p{0.30\textwidth}@{}}
\toprule
Claim & Status & Reason \\
\midrule
P versus NP & OPEN & unchanged \\
source-legal char-2 product-span rank & OPEN & not connected to a verified wall \\
compiler, arbitrary SAT decider to \((T,S,R)\) & OPEN & load-bearing \\
resolution: clause space \(\ge\) width minus \(k\) & THEOREM & Atserias--Dalmau, opened \\
resolution: short proofs are narrow & THEOREM & Ben-Sasson--Wigderson, opened \\
PC degree, immune / random \(k\)-CNF, char 2 & THEOREM & Alekhnovich--Razborov abstract, opened \\
PC space versus degree & NOT VERIFIED & no space paper opened \\
SAT not in \(n^{1+o(1)}\) time and \(n^{1-\varepsilon}\) space & THEOREM & Fortnow, opened \\
SAT time-space below the golden ratio & THEOREM & Fortnow--Lipton--van Melkebeek--Viglas \\
\(n^{o(1)}\)-space SAT time \(n^{2\cos(\pi/7)-o(1)}\) & THEOREM as cited & Williams manuscript, opened \\
alternation-trading cannot reach exponent 2 & METHOD LIMIT & cited there; 2012 paper not opened \\
\(BQF\) length \(1.0178n\) needs size \(\ge 2^{n/3}\) & THEOREM & Beame--Saks--Thathachar, opened \\
width meeting PC degree & CANDIDATE & no theorem opened \\
RAM space \(= \log(W_{\mathrm{BP}})\) for SAT & CANDIDATE & simulation not opened \\
\(D(n,m)\) in a general forbidden region & OPEN & compiler missing \\
semantic cut \(\ge 48\) on negative-star \(D(4,8)\) & RETIRED & exact cut 8 \\
\(Q_4\) as read-2 lower-bound candidate & RETIRED & \(\mathrm{lw}=7\), read-2 size \(\le 64\) \\
TEST-PROOFSPACE-D48 & UNRUN & no number produced \\
\bottomrule
\end{tabular}
}

\section*{Missing implications}
A feasible point of \(D(n,m)\) does not yield a forbidden point of SAT. A forbidden point of resolution does not yield a forbidden point of polynomial time. A small-space RAM lower bound does not yield a polynomial-space lower bound. An exponential branching-program size bound at length \(1.0178n\) for one quadratic form does not yield a SAT branching-program bound. Clause space meeting width does not compile contacts into a charged \((T,S,R)\) model.

\section*{Single next target}
Run TEST-PROOFSPACE-D48 on one frozen CNF, or prove there is no such CNF in the repository and stop. Do not treat an unrun limit as a separation.

% END_OF_VERIFIED_REPORT::GROK::2026-10-04_1149_PROOFSPACE1
\end{document}
